% TEMPLATE-MILC-v3.tex - plantilla maestra para todas las guias MILC v3
% Ver scripts/generadores/build-guia-milc.py para la lista completa de
% claves esperadas. El script valida que no queden placeholders sin
% reemplazar antes de escribir el archivo final.

\documentclass[11pt,letterpaper]{article}
\usepackage[margin=1.75cm]{geometry}
\usepackage[table]{xcolor}
\usepackage{fontspec}
\usepackage[spanish]{babel}
\usepackage{tikz}
\usepackage[most]{tcolorbox}
\usepackage{tabularx}
\usepackage{array}
\usepackage{fancyhdr}
\usepackage{titlesec}
\usepackage{ragged2e}
\usepackage{enumitem}
\usepackage{microtype}

\setmainfont{Helvetica Neue}
\setsansfont{Helvetica Neue}
\setlength{\parindent}{0pt}
\setlength{\parskip}{6pt}
\hyphenpenalty=200
\exhyphenpenalty=200
\pretolerance=400
\tolerance=2000
\setlength{\emergencystretch}{1em}
\renewcommand{\arraystretch}{1.25}
\setlist{leftmargin=5mm,itemsep=1mm,topsep=1mm}
\newcolumntype{Y}{>{\RaggedRight\arraybackslash}X}

\definecolor{milcMagenta}{HTML}{E6007E}
\definecolor{milcVino}{HTML}{5A0038}
\definecolor{milcTurquesa}{HTML}{0093A5}
\definecolor{milcVerde}{HTML}{58A923}
\definecolor{milcMostaza}{HTML}{E5B400}
\definecolor{milcOcre}{HTML}{B86600}
\definecolor{milcNegro}{HTML}{241816}
\definecolor{milcGris}{HTML}{F7F7F4}
\definecolor{milcLinea}{HTML}{E8E2DD}
\definecolor{milcGradePrimary}{HTML}{84CC16}
\definecolor{milcGradeDark}{HTML}{4D7C0F}
\definecolor{milcGradeSoft}{HTML}{ECFCCB}
\definecolor{milcGradeAccent}{HTML}{CFE7FF}
\definecolor{milcGradeText}{HTML}{FFFFFF}
\color{milcNegro}

\pagestyle{fancy}
\fancyhf{}
\lhead{\small Tecnología e Informática · Grado 7}
\rhead{\small Guía 3 · MILC v3}
\cfoot{\small\thepage}
\renewcommand{\headrulewidth}{0pt}

\newtcolorbox{titlebox}[1]{
  enhanced,colback=#1,colframe=#1,arc=7mm,boxrule=0pt,
  left=7mm,right=7mm,top=3mm,bottom=3mm,
  before skip=8pt,after skip=8pt,
  fontupper=\sffamily\bfseries\Large\color{white}
}

\newtcolorbox{softbox}[3]{
  enhanced,colback=#2,colframe=#1,borderline west={4pt}{0pt}{#1},
  arc=2mm,boxrule=.4pt,left=4mm,right=4mm,top=3mm,bottom=3mm,
  before skip=6pt,after skip=8pt,title={#3},coltitle=white,
  colbacktitle=#1,fonttitle=\sffamily\bfseries,fontupper=\RaggedRight,
  attach boxed title to top left={xshift=3mm,yshift=-2mm},
  boxed title style={arc=2mm, boxrule=0pt}
}

\newtcolorbox{drawbox}[2]{
  enhanced,colback=white,colframe=#1,arc=4mm,boxrule=1pt,
  left=5mm,right=5mm,top=4mm,bottom=4mm,before skip=8pt,after skip=8pt,
  title={#2},coltitle=white,colbacktitle=#1,fonttitle=\sffamily\bfseries,
  fontupper=\RaggedRight,
  attach boxed title to top left={xshift=4mm,yshift=-2mm},
  boxed title style={arc=3mm, boxrule=0pt}
}

\newtcolorbox{infoband}[1]{
  enhanced,colback=#1!8,colframe=#1!50,arc=2mm,boxrule=.5pt,
  left=4mm,right=4mm,top=2mm,bottom=2mm,before skip=4pt,after skip=4pt,
  fontupper=\small\RaggedRight
}

% Puente narrativo entre secciones (contrato editorial Conéctate).
% Da continuidad: "Acabas de X. Ahora vas a Y porque Z."
% Visualmente: banda izquierda en color del grado, texto en itálica pequeña.
\newtcolorbox{puentebox}{
  enhanced,colback=milcGradeSoft,colframe=milcGradeDark,
  borderline west={3pt}{0pt}{milcGradeDark},
  arc=0mm,boxrule=0pt,left=5mm,right=4mm,top=2.5mm,bottom=2.5mm,
  before skip=4pt,after skip=8pt,
  fontupper=\itshape\small\color{milcNegro}\RaggedRight
}
\newcommand{\puente}[1]{%
  \begin{puentebox}%
    \textcolor{milcGradeDark}{\textbf{\textrightarrow}}\hspace{2mm}#1%
  \end{puentebox}%
}

\newcommand{\linead}[1]{\noindent\color{milcLinea}\rule{#1}{0.5pt}\color{milcNegro}}
\newcommand{\respuesta}[1]{\par\vspace{2mm}\foreach \n in {1,...,#1}{\noindent\color{milcLinea}\rule{\linewidth}{0.5pt}\par\vspace{5mm}}\color{milcNegro}}
\newcommand{\checkbox}{\textcolor{milcGradeDark}{\fbox{\phantom{x}}}\hspace{5pt}}

\newcommand{\phasecard}[4]{
\begin{tcolorbox}[enhanced,colback=#3,colframe=#2,arc=3mm,boxrule=.5pt,height=3.05cm,valign=center,halign=center,left=1.5mm,right=1.5mm,top=2mm,bottom=2mm]
{\sffamily\bfseries\fontsize{22}{24}\selectfont\textcolor{#2}{#1}}\par
{\sffamily\bfseries\small #4}
\end{tcolorbox}
}

\begin{document}
\thispagestyle{empty}
\begin{tikzpicture}[remember picture,overlay]
  \fill[white] (current page.south west) rectangle (current page.north east);
  \path[fill=milcGradePrimary,rounded corners=16mm]
    ([xshift=.55cm,yshift=-.55cm]current page.north west) rectangle
    ([xshift=-.55cm,yshift=3.65cm]current page.south east);

  \node[anchor=north west,text=milcGradeText] at ([xshift=2.0cm,yshift=-1.85cm]current page.north west)
    {\sffamily\bfseries\fontsize{14}{16}\selectfont GRADO};
  \node[anchor=north west,text=milcGradeText,opacity=.82] at ([xshift=2.0cm,yshift=-2.75cm]current page.north west)
    {\sffamily\bfseries\fontsize{9}{11}\selectfont SERIE GUÍAS · TECNOLOGÍA E INFORMÁTICA};

  \begin{scope}[shift={([xshift=-3.05cm,yshift=-2.55cm]current page.north east)}]
    \fill[white,opacity=.80] (-.62,-.52) rectangle (.62,.52);
    \draw[milcGradeText,line width=1pt] (-.62,-.52) rectangle (.62,.52);
    \fill[milcGradeDark] (-.42,-.38) rectangle (-.22,.06);
    \fill[milcGradeAccent] (-.10,-.38) rectangle (.10,.24);
    \fill[milcGradePrimary!60!black] (.22,-.38) rectangle (.42,.38);
  \end{scope}

  \node[anchor=west,text=milcGradeText] at ([xshift=1.9cm,yshift=-12.85cm]current page.north west)
    {\sffamily\bfseries\fontsize{124}{124}\selectfont 7};
  \draw[milcGradeText,line width=1.7pt]
    ([xshift=4.52cm,yshift=-11.68cm]current page.north west) circle (.13cm);

  \node[anchor=west,text=milcGradeText,text width=8.7cm,align=left] at ([xshift=10.35cm,yshift=-11.6cm]current page.north west)
    {
     {\sffamily\bfseries\fontsize{19}{22}\selectfont OneDrive ---\\mi nube\\personal y\\cuarto de\\archivos}\\[4mm]
     {\sffamily\bfseries\fontsize{23}{26}\selectfont Guía 3-7-TIC}\\[2mm]
     {\sffamily\fontsize{13}{16}\selectfont Período 1 · Trabajo colaborativo en la nube}\\[5mm]
     {\sffamily\bfseries\fontsize{11}{13}\selectfont Institución Educativa Sor María Juliana}\\[1mm]
     {\sffamily\fontsize{10}{12}\selectfont Autor: Dr. Álvaro Cárdenas Orozco}
    };

  \path[fill=milcGradeSoft,rounded corners=7mm]
    ([xshift=1.35cm,yshift=.85cm]current page.south west) rectangle
    ([xshift=-1.35cm,yshift=4.25cm]current page.south east);
  \draw[milcGradeDark,line width=.65pt,rounded corners=7mm]
    ([xshift=1.35cm,yshift=.85cm]current page.south west) rectangle
    ([xshift=-1.35cm,yshift=4.25cm]current page.south east);
  \node[anchor=center,text=milcNegro,text width=17.0cm,align=center] at ([yshift=2.55cm]current page.south)
    {\sffamily\fontsize{10}{12}\selectfont Metodología MILC · Apertura ancestral · Pensamiento computacional · Triángulo Dussel-Estoicismo-Floridi\\
    \textcolor{milcGradeDark}{\bfseries Producto:} Tu \textbf{OneDrive activado} con cuenta institucional y \textbf{una estructura personal de 6 carpetas} creadas en la nube siguiendo el árbol que aprendiste en G6: \emph{Mis-documentos / Colegio / [materia] / [periodo]} más \emph{Personal / Imágenes / Descargas}. Más \textbf{2 archivos subidos a la nube} (cualquier Word o PDF que tengas) en su carpeta correcta. Más \textbf{en el cuaderno}: dibujo del árbol de OneDrive + lista de 5 ventajas de la nube frente al disco duro.};
\end{tikzpicture}
\mbox{}
\newpage

\section*{Apertura: OneDrive --- mi nube personal y mi cuarto de archivos}

\begin{softbox}{milcGradeDark}{milcGradeSoft}{Saber ancestral}
\textbf{Doña Carmen la costurera del barrio Cartago tenía un cuarto especial en su casa: el cuarto de las telas.} En una casa tradicional del Valle del Cauca, cuando uno entraba a su taller, lo primero que mostraba era \textbf{el cuarto de telas}: tela blanca para uniformes, tela floreada para vestidos, tela negra para sotanas, tela impermeable para gabanes. \emph{``Cada tela en su estante, cada estante con su etiqueta''}, decía doña Carmen. \textbf{¿Por qué tanta organización?} Porque cuando llegaba un cliente a las 7 de la mañana pidiendo \emph{``hágame un vestido de fiesta''}, ella iba derecho al estante de tela floreada, no perdía tiempo buscando. Si todo hubiera estado amontonado, habría pasado 30 minutos buscando antes de empezar a coser. \textbf{Hoy tu equivalente al cuarto de telas es OneDrive}: el espacio donde guardas tus archivos en orden. Si lo organizas bien desde el principio, encuentras todo en segundos. Si lo dejas amontonado, pierdes archivos importantes y vuelves a hacerlos.
\end{softbox}

\begin{softbox}{milcTurquesa}{milcGris}{Saber contemporáneo}
\textbf{OneDrive} es la \textbf{nube personal} que Microsoft te da con tu cuenta institucional. Es como un disco duro, pero \emph{en internet}: tus archivos viven en servidores de Microsoft, y los abres desde cualquier dispositivo con tu cuenta. \textbf{Te da 1 TB de espacio} (mil GB --- suficiente para guardar todos tus trabajos del colegio durante años). \textbf{5 ventajas claras sobre un disco duro local:} \emph{(1) accedes desde cualquier computador del mundo}; \emph{(2) si tu computador se daña, los archivos siguen vivos}; \emph{(3) coautoría --- varios pueden editar al mismo archivo}; \emph{(4) historial --- ves cada versión anterior y puedes restaurar}; \emph{(5) sincronización automática --- si lo abres en tu celular, ahí está}. \textbf{Lo que sigue siendo igual que en disco duro:} \emph{tú decides la estructura de carpetas}. Si la haces bien, OneDrive es un aliado. Si la haces mal, es un desorden mayor que el disco duro porque no lo ves físicamente.
\end{softbox}

\begin{softbox}{milcMostaza}{milcGris}{Pregunta-puente}
\textit{Si tu casa se incendia mañana y se quema tu computador, ¿\textbf{qué pasaría con tus tareas del año}? ¿Las has guardado solo en el disco local, o tienes copia en algún lugar?}
\end{softbox}

\begin{softbox}{milcVerde}{milcGris}{Saber y saber hacer}
Al terminar la clase: (1) podrás \textbf{identificar} las 5 ventajas de la nube; (2) sabrás \textbf{aplicar} las 5 reglas de nombrado (de G6·P2·S7) en OneDrive; (3) podrás \textbf{evaluar} si una estructura de carpetas está bien diseñada; (4) habrás \textbf{creado} tu OneDrive activo con 6 carpetas + 2 archivos.
\end{softbox}

\small
\begin{tabularx}{\textwidth}{>{\bfseries}p{3.0cm}Y}
\rowcolor{milcGradePrimary!12}Autor & Dr. Álvaro Cárdenas Orozco\\
DBA & Utiliza herramientas digitales para trabajar de manera colaborativa con otros (MEN, T\&I 7°)\\
Referentes & Saber ancestral de la minga · Microsoft 365 y OneDrive · Coautoría asincrónica\\
Producto final & Tu \textbf{OneDrive activado} con cuenta institucional y \textbf{una estructura personal de 6 carpetas} creadas en la nube siguiendo el árbol que aprendiste en G6: \emph{Mis-documentos / Colegio / [materia] / [periodo]} más \emph{Personal / Imágenes / Descargas}. Más \textbf{2 archivos subidos a la nube} (cualquier Word o PDF que tengas) en su carpeta correcta. Más \textbf{en el cuaderno}: dibujo del árbol de OneDrive + lista de 5 ventajas de la nube frente al disco duro.\\
\end{tabularx}
\normalsize

\puente{Hoy haces 4 cosas: comparas nube vs disco duro, activas OneDrive, creas la estructura, subes archivos.}

\begin{titlebox}{milcGradeDark}
Ruta MILC: así vamos a aprender
\end{titlebox}
\begin{tabularx}{\textwidth}{>{\centering\arraybackslash}X>{\centering\arraybackslash}X>{\centering\arraybackslash}X>{\centering\arraybackslash}X}
\phasecard{1}{milcVerde}{milcGris}{Escuta\\[-1mm]\small Observo, escucho y pregunto.} &
\phasecard{2}{milcTurquesa}{milcGris}{Sistematización\\[-1mm]\small Pensamiento computacional.} &
\phasecard{3}{milcMagenta}{milcGris}{Praxis\\[-1mm]\small Construyo y aplico.} &
\phasecard{4}{milcVino}{milcGris}{Evaluación\\[-1mm]\small Reflexión liberadora.}
\end{tabularx}


\puente{Empezamos con un caso real: ¿qué pasa si se daña tu computador?}

\begin{titlebox}{milcVerde}
Fase 1 · Escuta
\end{titlebox}

\begin{softbox}{milcVerde}{milcGris}{Escena para escuchar}
\textbf{Actividad 1 · ANALIZA --- 4 casos de pérdida de archivos} (12 min · individual). Lee estos 4 casos reales (cosas que pasan todos los días) y en el cuaderno responde: \emph{¿qué habría salvado a esa persona?}. \textbf{Caso 1:} A Lucía se le rompió el disco duro del computador la noche antes de entregar su trabajo final. Lo había hecho solo en local, no tenía copia en otro lado. \textbf{Caso 2:} A Pedro le robaron el computador. Junto al equipo se fue todo el archivo de fotos familiares de los últimos 5 años. \textbf{Caso 3:} A María se le cayó café sobre el portátil. Funciona el computador pero el disco no abre. \textbf{Caso 4:} A David se le perdió la memoria USB con la tesis de grado. Era la única copia.
\end{softbox}

\checkbox Leí los 4 casos con calma.\\
\checkbox Identifiqué qué habría salvado a cada persona.\\
\checkbox {'Pensé en tu propio escenario': '¿qué haría yo si me pasara?'}

\begin{infoband}{milcVerde}
\textbf{Lo que va al cuaderno (Actividad 1 · ANALIZA).} Título: «Actividad 1 --- 4 casos de pérdida». \textbf{Formato:} 4 respuestas cortas. \textbf{Extensión:} 4 líneas. \textbf{Sabes que terminaste cuando} reconoces el patrón común: la nube habría salvado a los 4.
\end{infoband}

\puente{Después aprendes qué es OneDrive y las 5 ventajas claves.}

\begin{titlebox}{milcTurquesa}
Fase 2 · Sistematización con pensamiento computacional
\end{titlebox}

\begin{softbox}{milcTurquesa}{milcGris}{Cuatro pilares aplicados al tema}
Los 4 casos comparten un mismo error: \textbf{no había copia en la nube}. Ahora aprendes qué es OneDrive y cómo te protege de esos escenarios.
\end{softbox}

\small
\begin{tabularx}{\textwidth}{>{\bfseries\RaggedRight\arraybackslash}p{3.6cm}Y}
\rowcolor{milcTurquesa!90}\color{white}Pilar & \color{white}\bfseries Aplicación al tema\\
1. Descomposición & \textbf{Qué es OneDrive técnicamente.} OneDrive es el \textbf{servicio de almacenamiento en la nube de Microsoft}. Cuando guardas un archivo en OneDrive, en realidad lo estás guardando en \textbf{servidores de Microsoft} (en centros de datos en Estados Unidos, Europa o Asia). Tu computador local mantiene solo una \emph{copia sincronizada}. Si pierdes el computador, los archivos siguen ahí. Si abres tu cuenta desde otro PC, ahí están. Es como tener un \textbf{cuarto de telas que existe simultáneamente en tu casa y en una bodega segura lejos}.\\
2. Reconocimiento de patrones & \textbf{Las 5 ventajas decisivas de OneDrive.} \textbf{(1) Acceso multidispositivo:} entras desde computador del colegio, computador de casa, celular, tableta. Mismos archivos siempre. \textbf{(2) Backup automático:} aunque borres tu disco duro, los archivos siguen vivos en la nube. \textbf{(3) Coautoría:} varias personas editan al mismo Word, Excel o PowerPoint a la vez (lo verás en S6). \textbf{(4) Historial de versiones:} OneDrive guarda automáticamente cada versión. Si dañas un documento, puedes restaurar la versión de ayer (lo verás en S8). \textbf{(5) Compartir fácil:} en lugar de mandar el archivo, mandas un enlace. El otro abre desde su computador sin descargar (lo verás en S5).\\
3. Abstracción & \textbf{Cómo activar OneDrive (paso a paso real).} \textbf{Opción A (en el navegador):} entra a \texttt{office.com}, inicia sesión con tu cuenta institucional (\texttt{tunombre@sormariajuliana.edu.co}). Aparece la página principal. Haz clic en el ícono de \emph{OneDrive} (nube azul). Ya estás dentro. \textbf{Opción B (instalando en el computador):} descarga \emph{OneDrive desktop} desde Microsoft. Te crea una carpeta especial en tu PC que se sincroniza automáticamente con la nube. Lo que pones en esa carpeta sube solo. \textbf{Recomendación inicial:} usa Opción A (navegador) este semestre. La B requiere más espacio en disco.\\
4. Algoritmo conceptual & \textbf{La estructura de carpetas para grado 7.} Aplica el principio del árbol jerárquico de G6·P2·S7. Estructura recomendada en OneDrive: \textbf{Mis-documentos} (raíz) → \emph{Colegio} → \emph{Tecnologia}, \emph{Matematicas}, \emph{Sociales} (cada materia con sus periodos) → \emph{Periodo-1}, \emph{Periodo-2}, \emph{Periodo-3}. Y carpetas paralelas: \emph{Personal}, \emph{Imagenes}, \emph{Descargas}. \textbf{Aplica las 5 reglas de nombrado:} sin espacios (usa guion), sin acentos, claro, con fecha si aplica, pensando en tu yo del próximo año.\\
\end{tabularx}
\normalsize

\begin{drawbox}{milcTurquesa}{5 ventajas + activación + estructura}
\textbf{5 ventajas:} acceso multidispositivo · backup auto · coautoría · historial · compartir fácil. \\[1mm]
\textbf{Activación:} office.com → iniciar sesión institucional → ícono OneDrive. \\[1mm]
\textbf{Estructura:} Mis-documentos → Colegio → Materias → Periodos.
\end{drawbox}

\begin{softbox}{milcMostaza}{milcGris}{Errores comunes y su corrección}
\textbf{Errores frecuentes al iniciar con OneDrive.} (1) \emph{No activar la cuenta el primer día} --- te toca hacerlo en momento de presión después. Hazlo ya. (2) \emph{Crear todas las carpetas al ras (sin jerarquía)} --- en un mes vas a tener 30 carpetas planas, sin orden. Mejor árbol desde el inicio. (3) \emph{Usar acentos y espacios} --- las URLs de los archivos compartidos se enredan. Sin acentos, sin espacios. (4) \emph{Subir archivos pesados sin pensar} --- ten en cuenta el espacio (1 TB suena mucho pero se llena con videos). (5) \emph{Olvidar la contraseña institucional} --- escríbela en lugar seguro. Sin contraseña, no entras.
\end{softbox}

\begin{infoband}{milcTurquesa}
\textbf{Lo que va al cuaderno (Actividad 2 · EXPLICA).} Título: «Actividad 2 --- OneDrive: las 5 ventajas». \textbf{Formato:} tabla de 5 filas, 2 columnas. Columna 1: nombre de la ventaja. Columna 2: ejemplo concreto de cómo te ayuda. \textbf{Extensión:} 5 filas. \textbf{Sabes que terminaste cuando} podrías convencer a tu mamá de empezar a usar OneDrive con 3 razones.
\end{infoband}

\puente{Con todo claro, abres OneDrive y armas tu estructura de carpetas personal.}

\begin{titlebox}{milcMagenta}
Fase 3 · Praxis con pensamiento computacional
\end{titlebox}

\begin{softbox}{milcMagenta}{milcGris}{Construyendo el producto con los cuatro pilares}
\textbf{Actividad 3 · APLICA --- OneDrive activo + 6 carpetas + 2 archivos} (30 min · individual). Vas a activar tu OneDrive (si no estaba) y crear la estructura de carpetas + subir 2 archivos. Si no tienes computador en clase, planificas en cuaderno y ejecutas en casa esta semana.
\end{softbox}

\small
\begin{tabularx}{\textwidth}{>{\bfseries\RaggedRight\arraybackslash}p{3.6cm}Y}
\rowcolor{milcMagenta!90}\color{white}Pilar & \color{white}\bfseries Aplicación al producto\\
Descomposición & \textbf{Paso 1 --- Entra a OneDrive.} Abre el navegador (Chrome, Edge). Escribe \texttt{office.com} en la barra. Inicia sesión con tu cuenta institucional (correo \texttt{...@sormariajuliana.edu.co} + contraseña). Si te pide cambiar la contraseña, hazlo (mínimo 8 caracteres, mayúsculas, números, signo --- igual a G6·P1·S3). Haz clic en el ícono de OneDrive (nube azul).\\
Patrones & \textbf{Paso 2 --- Crea la carpeta raíz \emph{Mis-documentos}.} Dentro de OneDrive, haz clic en \texttt{Nuevo → Carpeta}. Nómbrala \texttt{Mis-documentos} (sin acentos, sin espacios). Entra en ella.\\
Abstracción & \textbf{Paso 3 --- Crea las 4 subcarpetas (nivel 2).} Dentro de \texttt{Mis-documentos}, crea 4 carpetas: \texttt{Colegio}, \texttt{Personal}, \texttt{Imagenes}, \texttt{Descargas}. Estás en el nivel 2 del árbol.\\
Algoritmo de armado & \textbf{Paso 4 --- Subcarpetas de \emph{Colegio}.} Entra a \texttt{Colegio}. Crea: \texttt{Tecnologia} (sin acento, sin ñ). Entra. Adentro de \texttt{Tecnologia} crea: \texttt{Periodo-1}, \texttt{Periodo-2}, \texttt{Periodo-3}. Ya tienes el árbol completo de Tecnología. Puedes hacer lo mismo para otras materias si quieres (Matematicas, Sociales).\\
\end{tabularx}
\normalsize

\begin{drawbox}{milcMagenta}{Antes de cerrar la S3}
\checkbox Activé mi OneDrive entrando a office.com. \\[1mm]
\checkbox Creé la carpeta raíz Mis-documentos. \\[1mm]
\checkbox Creé 4 subcarpetas de nivel 2 (Colegio, Personal, etc.). \\[1mm]
\checkbox Dentro de Colegio creé Tecnologia con 3 periodos. \\[1mm]
\checkbox Subí al menos 2 archivos a una carpeta correcta. \\[1mm]
\checkbox Verifiqué desde otro dispositivo que se ven los archivos.
\end{drawbox}

\begin{softbox}{milcMagenta}{milcGris}{Plantilla de trabajo}
\textbf{Estructura objetivo.} \textit{Raíz:} Mis-documentos. \textit{Nivel 2:} Colegio + Personal + Imagenes + Descargas. \textit{Nivel 3 (dentro de Colegio):} Tecnologia con Periodo-1, Periodo-2, Periodo-3. \textit{Archivos:} 2 subidos a Tecnologia/Periodo-1.
\end{softbox}

\begin{infoband}{milcMagenta}
\textbf{Lo que va al cuaderno (Actividad 3 · APLICA).} Título: «Actividad 3 --- Mi OneDrive activo». \textbf{Formato:} dibujo del árbol + reflexión personal de 2 líneas + lista de archivos subidos. \textbf{Extensión:} 1 página. \textbf{Sabes que terminaste cuando} puedes abrir OneDrive desde cualquier dispositivo y ver tu estructura intacta.
\end{infoband}

\puente{El producto es OneDrive activo + 6 carpetas + 2 archivos + dibujo del árbol.}

\begin{titlebox}{milcMostaza}
Producto de la sesión
\end{titlebox}
\textbf{OneDrive activado con estructura de 6 carpetas + 2 archivos · firmado}.
\begin{softbox}{milcMostaza}{milcGris}{Criterios de calidad}
(1) OneDrive está activo con cuenta institucional. (2) Existe la carpeta raíz Mis-documentos. (3) Hay al menos 4 subcarpetas de nivel 2. (4) Dentro de Colegio existe Tecnologia con sus 3 periodos. (5) Hay 2+ archivos subidos a la carpeta correcta.
\end{softbox}

\puente{Antes de salir, abre tu OneDrive desde otro computador (o celular) para verificar que está sincronizado.}

\begin{titlebox}{milcVino}
Fase 4 · Evaluación liberadora — cinco dimensiones
\end{titlebox}

\begin{softbox}{milcVino}{milcGris}{Sentido de la evaluación}
Evaluar no es solo revisar una nota. Es reconocer qué aprendí, qué cambió en mí, qué emociones manejé, cómo me relaciono con la ciudad y la comunidad, y qué le dejaré a quien viene detrás.
\end{softbox}

\textbf{Mi reflexión liberadora en cinco dimensiones.} Completa cada frase iniciada:

\small
\begin{tabularx}{\textwidth}{>{\bfseries\RaggedRight\arraybackslash}p{3.4cm}Y>{\RaggedRight\arraybackslash}p{4.6cm}}
\rowcolor{milcVino}\color{white}Dimensión & \color{white}\bfseries Mi respuesta (frame starter) & \color{white}\bfseries Algo que descubrí\\
1. Personal & Lo que más cambió en mí fue \linead{6.5cm} & \linead{4.2cm}\\
2. Emocional & La emoción más fuerte fue \linead{2.5cm} y la regulé \linead{3cm} & \linead{4.2cm}\\
3. Ciudadana & En democracia esto importa porque \linead{6cm} & \linead{4.2cm}\\
4. Local (Cartago) & En mi barrio sería útil para \linead{6.5cm} & \linead{4.2cm}\\
5. Intergeneracional & A mi abuela/abuelo le diría \linead{6.5cm} & \linead{4.2cm}\\
\end{tabularx}
\normalsize

\begin{infoband}{milcVino}
\textbf{Para la próxima sustentación.} Vuelve a esta tabla en seis meses. Verás que las respuestas cambian — no porque lo aprendido cambie, sino porque tú cambias. La evaluación liberadora es una conversación contigo a través del tiempo.
\end{infoband}


\puente{Cerramos con 3 voces que te ayudan a entender por qué tener tu propio espacio digital es soberanía.}

\begin{titlebox}{milcGradeDark}
Triángulo de pensamiento · cierre filosófico
\end{titlebox}

\begin{softbox}{milcVino}{milcGris}{Dussel · lente del nosotros}
\textit{````Tener tu propio espacio --- físico o digital --- es la primera forma de soberanía.''''}

Hasta antes de OneDrive, tus archivos vivían \emph{en el computador del colegio} o \emph{en la USB que prestabas a un amigo} o \emph{en un Drive ajeno}. \textbf{No tenías un espacio propio}. Con OneDrive institucional, tienes tu \emph{cuarto de archivos digital}: 1 TB de espacio que es tuyo, organizado por ti, protegido por tu contraseña. Eso es soberanía digital --- el equivalente del cuarto propio que tenía doña Carmen en su casa.

\textbf{¿Estoy aprovechando mi espacio propio o sigo dejando mis archivos dispersos?}
\end{softbox}
\linead{\linewidth}\\[1mm]
\linead{\linewidth}

\begin{softbox}{milcMostaza}{milcGris}{Estoicismo · Marco Aurelio · cuidado interior}
\textit{````Lo que está ordenado por fuera ahorra el desorden interior. El espacio organizado da paz a la mente.''''}

Cuando entres a OneDrive y veas tu árbol de carpetas bien organizado, vas a sentir algo: \textbf{paz}. No tienes que pensar dónde guardar las cosas, dónde buscarlas, qué se perdió. \emph{El orden externo libera tu atención interna} para concentrarte en aprender. Es la diferencia entre estudiar en una mesa limpia y estudiar en una mesa con 50 cosas encima.

\textbf{¿Cuánta energía mental ahorraría si tuviera mi mundo digital ordenado?}
\end{softbox}
\linead{\linewidth}\\[1mm]
\linead{\linewidth}

\begin{softbox}{milcTurquesa}{milcGris}{Floridi · lente de la infoesfera}
\textit{````La nube no te da más espacio. Te da más libertad para llevar tu espacio a donde vayas.''''}

Antes, para acceder a tus archivos tenías que llevarte el computador o la USB. \textbf{Ahora con OneDrive llevas todo en tu cuenta}: vas al colegio, ahí están; vas a casa de la abuela, ahí están; te conectas desde la biblioteca, ahí están. \emph{Tu identidad digital se separa del dispositivo}. Eres tu cuenta + tu nube, no tu computador particular. Eso te hace móvil de una manera que tu abuela ni imaginaba.

\textbf{¿Cómo cambiaría mi vida si llevara mi mundo digital conmigo a todas partes?}
\end{softbox}
\linead{\linewidth}\\[1mm]
\linead{\linewidth}

\begin{titlebox}{milcVino}
Compromiso final
\end{titlebox}

\small
\begin{tabularx}{\textwidth}{>{\RaggedRight\arraybackslash}p{3.0cm}YYY}
\rowcolor{milcVino}\color{white}\bfseries Criterio & \color{white}\bfseries Lo logré & \color{white}\bfseries En proceso & \color{white}\bfseries Necesito apoyo\\
Comprensión del tema & Explico ideas centrales con mis palabras. & Reconozco ideas pero falta precisión. & Necesito repasar conceptos base.\\
Pensamiento computacional & Apliqué los 4 pilares al diseñar mi producto. & Apliqué algunos pilares. & Necesito releer la fase 2.\\
Aplicación y producto & Mi producto cumple los criterios y se puede revisar. & Producto funcional con ajustes pendientes. & Aún en construcción.\\
Reflexión 5 dimensiones & Respondí cada dimensión con honestidad. & Respondí algunas con detalle. & Mis respuestas son muy generales.\\
Triángulo de pensamiento & Conecté las 3 voces con mi trabajo real. & Conecté al menos una. & No relaciono las voces con mi proceso.\\
\end{tabularx}
\normalsize

\textbf{Hoy entendí que:}\\[1mm]
\linead{\linewidth}\\[1mm]
\linead{\linewidth}

\textbf{Mi compromiso para la próxima sesión es:}\\[1mm]
\linead{\linewidth}\\[1mm]
\linead{\linewidth}

\begin{softbox}{milcMostaza}{milcGris}{Cita de despedida · Marco Aurelio}
\textit{``El obstáculo en el camino se vuelve el camino. Lo que estorba al camino se vuelve camino.''}\\[1mm]
Cada error de hoy, bien observado, se convierte en pieza del camino que pisarás mañana.
\end{softbox}

\small
\begin{tabularx}{\textwidth}{>{\bfseries}p{3.6cm}Y}
\rowcolor{milcGradeDark!12}Nombre completo: & \linead{10cm}\\
Fecha: & \linead{4cm} \quad Curso/grupo: \linead{4cm}\\
Mi firma: & \linead{9cm}\\
\end{tabularx}
\normalsize

\begin{infoband}{milcGradeDark}
\textbf{Para la próxima sesión.} Esta semana, abre OneDrive al menos 3 veces desde dispositivos distintos (en el cole, en casa, en celular). Sube 5 archivos más a las carpetas correctas. La próxima clase aprendes a usar \textbf{Word, Excel y PowerPoint en línea}, directamente desde el navegador, guardando todo en OneDrive automáticamente.
\end{infoband}

\end{document}
